\chapter{État de l'Art et Cadre Théorique}
\section{Revue de littérature, ce qui se fait à l'Institut des Actuaires}

\subsection{Le socle traditionnel : les GLM et la Régression Logistique}
\subsection{Le changement de paradigme : l'hégémonie des modèles ensemblistes}
\subsection{L'ouverture vers l'apprentissage non supervisé}


\section{Les algorithmes d'IA pour la fraude}
\subsection{L'apprentissage supervisé : prédire la probabilité de fraude}

a. La Régression Logistique (Modèle de référence)

b. Les Arbres de décision et les Forêts Aléatoires (Random Forest)

c. Le Gradient Boosting (XGBoost)

\subsection{Apprentissage non supervisé : détection d'anomalies}
Détecter des anomalies / comportements atypiques (Isolation Forest, Clustering).
\subsection{Deep Learning \& Computer Vision}
Utilisation de YOLO pour l'analyse d'images de chocs et GANs pour la gestion des données déséquilibrées.

\section{Analye de réseaux (Graph Mining)}
\subsection{Pour détecter les bandes organisées}
Liens entre experts, garages, et assurés

